\chapter{Đặc Tả Yêu Cầu Phần Mềm}

% ============================================================
\section{Tổng Quan Sản Phẩm}
% ============================================================

\subsection{Mô Tả Hệ Thống}

\textbf{Tên hệ thống:} SmartDroneDelivery -- Nền tảng quản lý giao hàng bằng Drone tích hợp AI.

Hệ thống được xây dựng nhằm hỗ trợ doanh nghiệp logistics quản lý toàn bộ quy trình giao hàng
bằng máy bay không người lái, từ tiếp nhận yêu cầu giao hàng, quản lý khách hàng, kiện
hàng, trạm hạ cánh đến theo dõi trạng thái giao hàng theo thời gian thực và xác nhận hoàn tất.

\subsection{Các Thành Phần Hệ Thống}

Hệ thống gồm bốn thành phần chính:

\begin{itemize}
  \item \textbf{Website quản trị:} Dành cho Quản trị viên, Quản lý logistics,
        Điều phối viên và Nhân viên vận hành trạm. Cung cấp toàn bộ chức năng quản lý nghiệp vụ.
  \item \textbf{Giao diện Web cho người dùng cuối:} Dành cho Khách hàng tạo đơn, theo dõi hành trình và Nhân viên vận hành trạm xác nhận đơn. Tương thích responsive tốt trên trình duyệt của các thiết bị di động.
  \item \textbf{RESTful API:} Phục vụ trao đổi dữ liệu giữa giao diện web quản trị và giao diện web người dùng cuối.
        và các dịch vụ bên ngoài.
  \item \textbf{Dịch vụ AI:} Hỗ trợ ước lượng thời gian giao hàng dự kiến, chatbot trợ lý khách
        hàng và phân tích dữ liệu vận hành.
\end{itemize}

\subsection{Sơ Đồ Ngữ Cảnh Hệ Thống}

\begin{figure}[H]
  \centering
  \includegraphics[width=0.85\textwidth]{so-do/context_diagram.pdf}
  \caption{Sơ đồ ngữ cảnh hệ thống SmartDroneDelivery}
  \label{fig:context-diagram}
\end{figure}

% ============================================================
\section{Yêu Cầu Người Dùng}
% ============================================================

\subsection{Các Tác Nhân Hệ Thống}

Hệ thống SmartDroneDelivery có năm tác nhân chính tương tác với hệ thống:

\begin{table}[H]
  \centering
  \begin{tabularx}{\textwidth}{|l|l|>{\raggedright\arraybackslash}X|}
    \hline
    \textbf{Actor} & \textbf{Tên tiếng Việt} & \textbf{Mô tả vai trò} \\
    \hline
    Admin             & Quản trị viên              & Quản lý hệ thống, cấu hình, phân quyền người dùng và giám sát hệ thống \\ \hline
    Logistics Manager & Quản lý logistics          & Theo dõi báo cáo, thống kê, giám sát hiệu suất và đánh giá hoạt động giao hàng \\ \hline
    Dispatcher        & Điều phối viên             & Tiếp nhận, phê duyệt, lên lịch và điều phối đơn giao hàng; xử lý sự cố giao hàng thất bại \\ \hline
    Station Operator  & Nhân viên vận hành trạm    & Xác nhận kiện hàng tại trạm, cập nhật trạng thái và hỗ trợ xác nhận giao hàng thành công \\ \hline
    Customer          & Khách hàng                 & Đăng ký tài khoản, tạo đơn giao hàng, theo dõi tiến trình và xác nhận khi nhận hàng \\
    \hline
  \end{tabularx}
  \caption{Bảng các tác nhân hệ thống}
  \label{tab:actors}
\end{table}

\subsection{Sơ Đồ Use Case}
\subsubsection{Use Case -- Quản Lý Khách Hàng Và Đơn Hàng Giao}

\begin{figure}[H]
  \centering
  \includegraphics[width=0.88\textwidth]{so-do/uc_customer_order.pdf}
  \caption{Sơ đồ use case quản lý khách hàng và đơn hàng giao}
  \label{fig:uc-customer-order}
\end{figure}

\subsubsection{Use Case -- Quản Lý Kiện Hàng}

\begin{figure}[H]
  \centering
  \includegraphics[width=0.88\textwidth]{so-do/uc_package_management.pdf}
  \caption{Sơ đồ use case quản lý kiện hàng}
  \label{fig:uc-package}
\end{figure}

\subsubsection{Use Case -- Quản Lý Trạm Hạ Cánh}

\begin{figure}[H]
  \centering
  \includegraphics[width=0.88\textwidth]{so-do/uc_station_management.pdf}
  \caption{Sơ đồ use case quản lý trạm hạ cánh}
  \label{fig:uc-station}
\end{figure}

\subsubsection{Use Case -- Theo Dõi Và Xác Nhận Giao Hàng}

\begin{figure}[H]
  \centering
  \includegraphics[width=0.88\textwidth]{so-do/uc_tracking_confirmation.pdf}
  \caption{Sơ đồ use case theo dõi và xác nhận giao hàng}
  \label{fig:uc-tracking}
\end{figure}

\subsubsection{Use Case -- Phân Tích Và Dịch Vụ AI}

\begin{figure}[H]
  \centering
  \includegraphics[width=0.88\textwidth]{so-do/uc_ai_analytics.pdf}
  \caption{Sơ đồ use case phân tích và dịch vụ AI}
  \label{fig:uc-ai}
\end{figure}

\subsection{Danh Sách Use Case Tổng Hợp}

\begin{table}[H]
  \centering
  \begin{tabularx}{\textwidth}{|l|X|l|}
    \hline
    \textbf{UC ID} & \textbf{Tên Use Case} & \textbf{Actor chính} \\
    \hline
    UC-01 & Đăng nhập và Phân quyền                  & Tất cả \\ \hline
    UC-02 & Tạo đơn giao hàng mới                    & Customer \\ \hline
    UC-03 & Cập nhật và Hủy đơn giao hàng             & Customer \\ \hline
    UC-04 & Phê duyệt yêu cầu giao hàng              & Dispatcher \\ \hline
    UC-05 & Lập lịch giao hàng                       & Dispatcher \\ \hline
    UC-06 & Xử lý giao hàng thất bại                 & Dispatcher \\ \hline
    UC-07 & Xác nhận nhận/gửi kiện hàng tại trạm     & Station Operator \\ \hline
    UC-08 & Theo dõi giao hàng thời gian thực        & Customer, Dispatcher, Station Operator \\ \hline
    UC-09 & Xác nhận giao hàng thành công            & Station Operator, Customer \\ \hline
    UC-10 & Ước lượng thời gian giao hàng dự kiến    & Customer, Dispatcher \\ \hline
    UC-11 & Trợ lý AI cho khách hàng                 & Customer \\ \hline
    UC-12 & Tạo và quản lý tài khoản người dùng      & Admin \\ \hline
    UC-13 & Quản lý thiết bị Drone                   & Admin \\ \hline
    UC-14 & Quản lý trạm hạ cánh                     & Admin, Logistics Manager \\ \hline
    UC-15 & Xem báo cáo và thống kê      & Logistics Manager, Admin \\ \hline
    UC-16 & Quản lý địa chỉ giao hàng                & Customer \\ \hline
    UC-17 & Đăng ký tài khoản mới                    & Customer \\ \hline
    UC-18 & Quản lý thông tin cá nhân                & Tất cả \\ \hline
    UC-19 & Tóm tắt thông tin giao hàng bằng AI      & Customer, Dispatcher, Logistics Manager \\
    \hline
  \end{tabularx}
  \caption{Bảng tổng hợp use case hệ thống SmartDroneDelivery}
  \label{tab:usecases}
\end{table}

% ============================================================
\section{Đặc Tả Use Case}
% ============================================================

% --- UC-01 ---
\subsection{UC-01: Đăng Nhập và Phân Quyền}
\begin{table}[H]
  \centering
  \begin{tabularx}{\textwidth}{|l|X|}
    \hline
    \multicolumn{2}{|l|}{\textbf{Đăng Nhập và Phân Quyền}} \\
    \hline
    \textbf{Tác nhân}           & Tất cả tác nhân như Khách hàng, Điều phối viên, Nhân viên vận hành trạm, Quản lý logistics và Quản trị viên \\ \hline
    \textbf{Mô tả}              & Người dùng đăng nhập vào hệ thống bằng tài khoản đã được cấp; hệ thống xác định vai trò và phân quyền tương ứng. \\ \hline
    \textbf{Điều kiện tiên quyết} & Người dùng đã có tài khoản hợp lệ trong hệ thống. \\ \hline
    \textbf{Luồng sự kiện chính} &
      \begin{enumerate}[leftmargin=*, nosep, topsep=0pt]
        \item Người dùng nhập email/tên đăng nhập và mật khẩu.
        \item Hệ thống xác thực thông tin đăng nhập.
        \item Hệ thống xác định vai trò của tài khoản.
        \item Hệ thống chuyển hướng đến giao diện tương ứng với vai trò.
      \end{enumerate} \\ \hline
    \textbf{Luồng rẽ nhánh / Ngoại lệ} &
      \textbf{2a.} Sai tài khoản hoặc mật khẩu $\rightarrow$ hệ thống báo lỗi, yêu cầu nhập lại.
      \newline
      \textbf{2b.} Tài khoản bị khóa $\rightarrow$ hệ thống thông báo không thể đăng nhập. \\ \hline
    \textbf{Hậu điều kiện}      & Người dùng được đăng nhập thành công, có phiên làm việc hợp lệ. \\
    \hline
  \end{tabularx}
  \caption{Đặc tả Đăng Nhập và Phân Quyền}
  \label{tab:uc01}
\end{table}

% --- UC-02 ---
\subsection{UC-02: Tạo Đơn Giao Hàng Mới}
\begin{table}[H]
  \centering
  \begin{tabularx}{\textwidth}{|l|X|}
    \hline
    \multicolumn{2}{|l|}{\textbf{Tạo Đơn Giao Hàng Mới}} \\
    \hline
    \textbf{Tác nhân}           & Khách hàng \\ \hline
    \textbf{Mô tả}              & Khách hàng tạo một yêu cầu giao hàng mới trên hệ thống. \\ \hline
    \textbf{Điều kiện tiên quyết} & Đã đăng nhập; đã có ít nhất 1 địa chỉ giao hàng. \\ \hline
    \textbf{Luồng sự kiện chính} &
      \begin{enumerate}[leftmargin=*, nosep, topsep=0pt]
        \item Khách hàng chọn ``Tạo đơn hàng mới''.
        \item Nhập địa chỉ giao hàng, mô tả gói hàng, cân nặng, thời gian mong muốn.
        \item Hệ thống kiểm tra cân nặng gói hàng hợp lệ.
        \item Hệ thống tạo đơn hàng với trạng thái = ``Chờ duyệt''.
        \item Hệ thống thông báo cho Dispatcher về đơn hàng mới.
      \end{enumerate} \\ \hline
    \textbf{Luồng rẽ nhánh / Ngoại lệ} &
      \textbf{3a.} Gói hàng vượt tải trọng cho phép $\rightarrow$ báo lỗi, không tạo đơn.
      \newline
      \textbf{3b.} Địa chỉ ngoài vùng phục vụ $\rightarrow$ báo lỗi, không tạo đơn. \\ \hline
    \textbf{Hậu điều kiện}      & Đơn hàng mới được lưu vào hệ thống với trạng thái ``Chờ duyệt''. \\
    \hline
  \end{tabularx}
  \caption{Đặc tả Tạo Đơn Giao Hàng Mới}
  \label{tab:uc02}
\end{table}

% --- UC-03 ---
\subsection{UC-03: Cập Nhật hoặc Hủy Đơn Giao Hàng}
\begin{table}[H]
  \centering
  \begin{tabularx}{\textwidth}{|l|X|}
    \hline
    \multicolumn{2}{|l|}{\textbf{Cập Nhật hoặc Hủy Đơn Giao Hàng}} \\
    \hline
    \textbf{Tác nhân}           & Khách hàng \\ \hline
    \textbf{Mô tả}              & Khách hàng sửa thông tin hoặc hủy đơn khi đơn chưa được xử lý. \\ \hline
    \textbf{Điều kiện tiên quyết} & Đơn đang ở trạng thái ``Chờ duyệt'' khi chưa gán trạm hoặc drone. \\ \hline
    \textbf{Luồng sự kiện chính} &
      \begin{enumerate}[leftmargin=*, nosep, topsep=0pt]
        \item Khách hàng mở chi tiết đơn hàng.
        \item Chọn ``Sửa'' hoặc ``Hủy''.
        \item Nếu sửa: cập nhật lại thông tin, lưu thay đổi.
        \item Nếu hủy: xác nhận, hệ thống đổi trạng thái = ``Đã hủy''.
      \end{enumerate} \\ \hline
    \textbf{Luồng rẽ nhánh / Ngoại lệ} &
      \textbf{2a.} Đơn đã được duyệt hoặc gán trạm $\rightarrow$ không cho sửa hoặc hủy, báo ``Đơn đã xử lý, không thể thay đổi''. \\ \hline
    \textbf{Hậu điều kiện}      & Đơn được cập nhật hoặc chuyển trạng thái ``Đã hủy''. \\
    \hline
  \end{tabularx}
  \caption{Đặc tả Cập Nhật hoặc Hủy Đơn Giao Hàng}
  \label{tab:uc03}
\end{table}

% --- UC-04 ---
\subsection{UC-04: Phê Duyệt Yêu Cầu Giao Hàng}
\begin{table}[H]
  \centering
  \begin{tabularx}{\textwidth}{|l|X|}
    \hline
    \multicolumn{2}{|l|}{\textbf{Phê Duyệt Yêu Cầu Giao Hàng}} \\
    \hline
    \textbf{Tác nhân}           & Điều phối viên \\ \hline
    \textbf{Mô tả}              & Dispatcher xem xét và duyệt hoặc từ chối đơn hàng mới. \\ \hline
    \textbf{Điều kiện tiên quyết} & Đơn đang ở trạng thái ``Chờ duyệt''. \\ \hline
    \textbf{Luồng sự kiện chính} &
      \begin{enumerate}[leftmargin=*, nosep, topsep=0pt]
        \item Dispatcher xem danh sách đơn ``Chờ duyệt''.
        \item Mở chi tiết đơn, kiểm tra thông tin gói hàng, địa chỉ.
        \item Bấm ``Duyệt'' $\rightarrow$ trạng thái chuyển ``Đã duyệt''.
        \item Hoặc bấm ``Từ chối'' $\rightarrow$ nhập lý do, trạng thái = ``Bị từ chối''.
      \end{enumerate} \\ \hline
    \textbf{Luồng rẽ nhánh / Ngoại lệ} &
      \textbf{3a.} Không có trạm hạ cánh khả dụng gần địa chỉ giao $\rightarrow$ hệ thống cảnh báo trước khi duyệt. \\ \hline
    \textbf{Hậu điều kiện}      & Đơn chuyển trạng thái ``Đã duyệt'' hoặc ``Bị từ chối''. \\
    \hline
  \end{tabularx}
  \caption{Đặc tả Phê Duyệt Yêu Cầu Giao Hàng}
  \label{tab:uc04}
\end{table}

% --- UC-05 ---
\subsection{UC-05: Lập Lịch Giao Hàng}
\begin{table}[H]
  \centering
  \begin{tabularx}{\textwidth}{|l|X|}
    \hline
    \multicolumn{2}{|l|}{\textbf{Lập Lịch Giao Hàng}} \\
    \hline
    \textbf{Tác nhân}           & Điều phối viên \\ \hline
    \textbf{Mô tả}              & Dispatcher xếp lịch giao hàng cho đơn đã duyệt, gán trạm và thời gian dự kiến. \\ \hline
    \textbf{Điều kiện tiên quyết} & Đơn đang ở trạng thái ``Đã duyệt''. \\ \hline
    \textbf{Luồng sự kiện chính} &
      \begin{enumerate}[leftmargin=*, nosep, topsep=0pt]
        \item Dispatcher chọn đơn đã duyệt.
        \item Hệ thống gợi ý trạm hạ cánh gần địa chỉ giao và đang hoạt động.
        \item Dispatcher chọn trạm và thời gian giao dự kiến.
        \item Hệ thống cập nhật trạng thái đơn = ``Đã lên lịch''.
      \end{enumerate} \\ \hline
    \textbf{Luồng rẽ nhánh / Ngoại lệ} &
      \textbf{2a.} Không có trạm khả dụng $\rightarrow$ hệ thống báo, Dispatcher phải chọn trạm khác hoặc hoãn lịch. \\ \hline
    \textbf{Hậu điều kiện}      & Đơn được gán trạm và thời gian giao, trạng thái ``Đã lên lịch''. \\
    \hline
  \end{tabularx}
  \caption{Đặc tả Lập Lịch Giao Hàng}
  \label{tab:uc05}
\end{table}

% --- UC-06 ---
\subsection{UC-06: Xử Lý Giao Hàng Thất Bại}
\begin{table}[H]
  \centering
  \begin{tabularx}{\textwidth}{|l|X|}
    \hline
    \multicolumn{2}{|l|}{\textbf{Xử Lý Giao Hàng Thất Bại}} \\
    \hline
    \textbf{Tác nhân}           & Điều phối viên \\ \hline
    \textbf{Mô tả}              & Dispatcher xử lý các đơn hàng giao không thành công, được mở rộng từ UC-08 Theo dõi tiến trình giao hàng. \\ \hline
    \textbf{Điều kiện tiên quyết} & Đơn đang ở trạng thái ``Đang giao'' và phát sinh sự cố. \\ \hline
    \textbf{Luồng sự kiện chính} &
      \begin{enumerate}[leftmargin=*, nosep, topsep=0pt]
        \item Hệ thống hoặc Nhân viên vận hành trạm ghi nhận sự cố giao hàng.
        \item Dispatcher xem chi tiết sự cố.
        \item Dispatcher quyết định: giao lại, hủy đơn, hoặc liên hệ khách hàng.
        \item Hệ thống cập nhật trạng thái đơn tương ứng.
      \end{enumerate} \\ \hline
    \textbf{Luồng rẽ nhánh / Ngoại lệ} &
      \textbf{3a.} Sự cố nghiêm trọng (mất kiện hàng) $\rightarrow$ chuyển cho Logistics Manager xử lý thêm. \\ \hline
    \textbf{Hậu điều kiện}      & Đơn được xử lý lại hoặc chuyển trạng thái ``Giao thất bại''. \\
    \hline
  \end{tabularx}
  \caption{Đặc tả Xử Lý Giao Hàng Thất Bại}
  \label{tab:uc06}
\end{table}

% --- UC-07 ---
\subsection{UC-07: Xác Nhận Nhận/Gửi Kiện Hàng Tại Trạm}
\begin{table}[H]
  \centering
  \begin{tabularx}{\textwidth}{|l|X|}
    \hline
    \multicolumn{2}{|l|}{\textbf{Xác Nhận Nhận và Gửi Kiện Hàng Tại Trạm}} \\
    \hline
    \textbf{Tác nhân}           & Nhân viên vận hành trạm \\ \hline
    \textbf{Mô tả}              & Nhân viên trạm xác nhận đã nhận kiện hàng từ khách hoặc đặt kiện hàng lên drone để giao. \\ \hline
    \textbf{Điều kiện tiên quyết} & Đơn đã được lên lịch, gán đúng trạm này. \\ \hline
    \textbf{Luồng sự kiện chính} &
      \begin{enumerate}[leftmargin=*, nosep, topsep=0pt]
        \item Nhân viên trạm nhận kiện hàng từ khách hàng hoặc từ kho.
        \item Quét hoặc nhập mã đơn hàng để xác nhận.
        \item Hệ thống cập nhật trạng thái kiện hàng = ``Đã nhận tại trạm''.
        \item Khi drone sẵn sàng, nhân viên xác nhận gửi đi, trạng thái = ``Đang giao''.
      \end{enumerate} \\ \hline
    \textbf{Luồng rẽ nhánh / Ngoại lệ} &
      \textbf{2a.} Kiện hàng không khớp mã đơn $\rightarrow$ báo lỗi, không xác nhận. \\ \hline
    \textbf{Hậu điều kiện}      & Trạng thái kiện hàng được cập nhật đúng theo bước xử lý tại trạm. \\
    \hline
  \end{tabularx}
  \caption{Đặc tả Xác Nhận Nhận và Gửi Kiện Hàng Tại Trạm}
  \label{tab:uc07}
\end{table}

% --- UC-08 ---
\subsection{UC-08: Theo Dõi Giao Hàng Thời Gian Thực}
\begin{table}[H]
  \centering
  \begin{tabularx}{\textwidth}{|l|X|}
    \hline
    \multicolumn{2}{|l|}{\textbf{Theo Dõi Giao Hàng Thời Gian Thực}} \\
    \hline
    \textbf{Tác nhân}           & Khách hàng, Điều phối viên, Nhân viên vận hành trạm \\ \hline
    \textbf{Mô tả}              & Xem vị trí và trạng thái di chuyển của đơn hàng đang giao, cập nhật theo thời gian thực. \\ \hline
    \textbf{Điều kiện tiên quyết} & Đơn đang ở trạng thái ``Đang giao''. \\ \hline
    \textbf{Luồng sự kiện chính} &
      \begin{enumerate}[leftmargin=*, nosep, topsep=0pt]
        \item Actor mở màn hình theo dõi đơn hàng.
        \item Hệ thống lấy tọa độ mới nhất qua Supabase Realtime.
        \item Bản đồ cập nhật vị trí di chuyển liên tục.
        \item Khi đơn hoàn tất, dừng cập nhật vị trí.
      \end{enumerate} \\ \hline
    \textbf{Luồng rẽ nhánh / Ngoại lệ} &
      \textbf{2a.} Mất kết nối dữ liệu vị trí $\rightarrow$ hiển thị trạng thái cuối cùng ghi nhận được, báo ``Đang chờ cập nhật''. \\ \hline
    \textbf{Hậu điều kiện}      & Actor xem được vị trí hoặc trạng thái giao hàng mới nhất. \\
    \hline
  \end{tabularx}
  \caption{Đặc tả Theo Dõi Giao Hàng Thời Gian Thực}
  \label{tab:uc08}
\end{table}

% --- UC-09 ---
\subsection{UC-09: Xác Nhận Giao Hàng Thành Công}
\begin{table}[H]
  \centering
  \begin{tabularx}{\textwidth}{|l|X|}
    \hline
    \multicolumn{2}{|l|}{\textbf{Xác Nhận Giao Hàng Thành Công}} \\
    \hline
    \textbf{Tác nhân}           & Nhân viên vận hành trạm, Khách hàng \\ \hline
    \textbf{Mô tả}              & Xác nhận đơn hàng đã được giao đến tay người nhận. \\ \hline
    \textbf{Điều kiện tiên quyết} & Drone đã đến điểm giao khi đơn ở trạng thái ``Đang giao'' và gần hoàn tất. \\ \hline
    \textbf{Luồng sự kiện chính} &
      \begin{enumerate}[leftmargin=*, nosep, topsep=0pt]
        \item Nhân viên trạm hoặc Khách hàng xác nhận đã nhận được kiện hàng.
        \item Hệ thống cập nhật trạng thái đơn = ``Đã giao thành công''.
        \item Hệ thống gửi thông báo hoàn tất cho Khách hàng.
      \end{enumerate} \\ \hline
    \textbf{Luồng rẽ nhánh / Ngoại lệ} &
      \textbf{1a.} Không xác nhận trong thời gian quy định $\rightarrow$ hệ thống tự động nhắc nhở hoặc chuyển trạng thái ``Chờ xác nhận quá hạn''. \\ \hline
    \textbf{Hậu điều kiện}      & Đơn hàng ở trạng thái ``Đã giao thành công'', quy trình giao hàng kết thúc. \\
    \hline
  \end{tabularx}
  \caption{Đặc tả Xác Nhận Giao Hàng Thành Công}
  \label{tab:uc09}
\end{table}

% --- UC-10 ---
\subsection{UC-10: Ước Lượng Thời Gian Giao Hàng Dự Kiến}
\begin{table}[H]
  \centering
  \begin{tabularx}{\textwidth}{|l|X|}
    \hline
    \multicolumn{2}{|l|}{\textbf{Ước Lượng Thời Gian Giao Hàng}} \\
    \hline
    \textbf{Tác nhân}           & Khách hàng, Điều phối viên \\ \hline
    \textbf{Mô tả}              & Hệ thống ước tính thời gian dự kiến giao hàng dựa trên khoảng cách và tốc độ trung bình. \\ \hline
    \textbf{Điều kiện tiên quyết} & Đơn đã được lên lịch, có tọa độ trạm xuất phát và địa chỉ giao. \\ \hline
    \textbf{Luồng sự kiện chính} &
      \begin{enumerate}[leftmargin=*, nosep, topsep=0pt]
        \item Hệ thống lấy tọa độ trạm và địa chỉ giao.
        \item Tính khoảng cách và áp dụng công thức ước lượng tốc độ trung bình.
        \item Hiển thị ETA cho Khách hàng/Dispatcher.
        \item ETA được cập nhật lại khi có thay đổi vị trí thực tế.
      \end{enumerate} \\ \hline
    \textbf{Luồng rẽ nhánh / Ngoại lệ} &
      \textbf{1a.} Thiếu dữ liệu tọa độ $\rightarrow$ hiển thị ``Chưa thể ước lượng''. \\ \hline
    \textbf{Hậu điều kiện}      & Actor xem được thời gian dự kiến giao hàng. \\
    \hline
  \end{tabularx}
  \caption{Đặc tả Ước Lượng Thời Gian Giao Hàng}
  \label{tab:uc10}
\end{table}

% --- UC-11 ---
\subsection{UC-11: Trợ Lý AI Cho Khách Hàng}
\begin{table}[H]
  \centering
  \begin{tabularx}{\textwidth}{|l|X|}
    \hline
    \multicolumn{2}{|l|}{\textbf{Trợ Lý AI Cho Khách Hàng}} \\
    \hline
    \textbf{Tác nhân}           & Khách hàng \\ \hline
    \textbf{Mô tả}              & Khách hàng đặt câu hỏi liên quan đến đơn hàng; chatbot dựa trên LLM API trả lời tự động bao gồm tra cứu thông tin giao hàng. \\ \hline
    \textbf{Điều kiện tiên quyết} & Khách hàng đã đăng nhập. \\ \hline
    \textbf{Luồng sự kiện chính} &
      \begin{enumerate}[leftmargin=*, nosep, topsep=0pt]
        \item Khách hàng mở khung chat, nhập câu hỏi.
        \item Hệ thống gửi câu hỏi kèm ngữ cảnh đơn hàng đến LLM API.
        \item Hệ thống tra cứu thông tin đơn hàng liên quan nếu cần.
        \item Chatbot trả lời câu hỏi của khách hàng.
      \end{enumerate} \\ \hline
    \textbf{Luồng rẽ nhánh / Ngoại lệ} &
      \textbf{2a.} LLM API không phản hồi/lỗi $\rightarrow$ hiển thị ``Chatbot tạm thời không khả dụng'', gợi ý liên hệ Dispatcher. \\ \hline
    \textbf{Hậu điều kiện}      & Khách hàng nhận được câu trả lời liên quan đến đơn hàng. \\
    \hline
  \end{tabularx}
  \caption{Đặc tả Trợ Lý AI Cho Khách Hàng}
  \label{tab:uc11}
\end{table}

% --- UC-12 ---
\subsection{UC-12: Tạo và Quản Lý Tài Khoản Người Dùng}
\begin{table}[H]
  \centering
  \begin{tabularx}{\textwidth}{|l|X|}
    \hline
    \multicolumn{2}{|l|}{\textbf{Tạo và Quản Lý Tài Khoản Người Dùng}} \\
    \hline
    \textbf{Tác nhân}           & Quản trị viên hệ thống \\ \hline
    \textbf{Mô tả}              & Admin tạo, chỉnh sửa, phân quyền tài khoản cho các actor trong hệ thống. \\ \hline
    \textbf{Điều kiện tiên quyết} & Admin đã đăng nhập với quyền quản trị. \\ \hline
    \textbf{Luồng sự kiện chính} &
      \begin{enumerate}[leftmargin=*, nosep, topsep=0pt]
        \item Admin mở màn hình quản lý người dùng.
        \item Tạo tài khoản mới, gán vai trò như Khách hàng, Điều phối viên, Nhân viên vận hành trạm, Quản lý logistics hoặc Quản trị viên.
        \item Hoặc chỉnh sửa hoặc khóa tài khoản hiện có.
        \item Hệ thống lưu thay đổi và ghi log hoạt động.
      \end{enumerate} \\ \hline
    \textbf{Luồng rẽ nhánh / Ngoại lệ} &
      \textbf{2a.} Email hoặc tài khoản đã tồn tại $\rightarrow$ báo lỗi, yêu cầu nhập lại. \\ \hline
    \textbf{Hậu điều kiện}      & Tài khoản được tạo/cập nhật, có hiệu lực ngay trong hệ thống. \\
    \hline
  \end{tabularx}
  \caption{Đặc tả Tạo và Quản Lý Tài Khoản Người Dùng}
  \label{tab:uc12}
\end{table}

% --- UC-13 ---
\subsection{UC-13: Quản lý thiết bị Drone}
\begin{table}[H]
  \centering
  \begin{tabularx}{\textwidth}{|l|X|}
    \hline
    \multicolumn{2}{|l|}{\textbf{Quản lý thiết bị Drone}} \\
    \hline
    \textbf{Tác nhân}           & Quản trị viên \\ \hline
    \textbf{Mô tả}              & Cho phép Admin thêm mới, chỉnh sửa thông tin, xóa hoặc cập nhật trạng thái bảo trì/hỏng hóc của thiết bị drone. \\ \hline
    \textbf{Điều kiện tiên quyết} & Admin đã đăng nhập thành công vào hệ thống. \\ \hline
    \textbf{Luồng sự kiện chính} &
      \begin{enumerate}[leftmargin=*, nosep, topsep=0pt]
        \item Admin mở giao diện quản lý danh mục Drone.
        \item Hệ thống hiển thị danh sách các drone kèm theo thông tin trạng thái, dung lượng pin, ngày bảo trì gần nhất.
        \item Nếu Admin muốn thêm drone mới: nhập các thông số ban đầu $\rightarrow$ Hệ thống lưu vào cơ sở dữ liệu.
        \item Nếu Admin muốn sửa thông tin: chọn thiết bị, cập nhật trạng thái hoạt động như Sẵn sàng, Bảo trì hoặc Hỏng hoặc cập nhật dung lượng pin $\rightarrow$ Hệ thống lưu thay đổi.
        \item Nếu Admin muốn xóa: chọn thiết bị $\rightarrow$ Hệ thống xóa bản ghi hoặc đổi trạng thái ngừng hoạt động.
      \end{enumerate} \\ \hline
    \textbf{Luồng rẽ nhánh / Ngoại lệ} &
      \textbf{3a/4a.} Trạng thái hoặc dung lượng pin nhập không hợp lệ, ví dụ pin ngoài khoảng từ 0 đến 100 $\rightarrow$ Hệ thống báo lỗi, yêu cầu chỉnh sửa lại.
      \newline
      \textbf{5a.} Drone đang thực hiện chuyến bay giao hàng ở trạng thái ``Đang giao'' $\rightarrow$ Hệ thống cảnh báo không được xóa hoặc đổi sang trạng thái bảo trì/hỏng. \\ \hline
    \textbf{Hậu điều kiện}      & Thông tin thiết bị drone được cập nhật thành công vào cơ sở dữ liệu. \\
    \hline
  \end{tabularx}
  \caption{Đặc tả Quản lý thiết bị Drone}
  \label{tab:uc13}
\end{table}

% --- UC-14 ---
\subsection{UC-14: Quản lý trạm hạ cánh}
\begin{table}[H]
  \centering
  \begin{tabularx}{\textwidth}{|l|X|}
    \hline
    \multicolumn{2}{|l|}{\textbf{Quản lý trạm hạ cánh}} \\
    \hline
    \textbf{Tác nhân}           & Quản trị viên, Quản lý logistics \\ \hline
    \textbf{Mô tả}              & Cho phép Admin hoặc Logistics Manager thêm mới, chỉnh sửa thông tin hoặc thay đổi trạng thái hoạt động của trạm hạ cánh. \\ \hline
    \textbf{Điều kiện tiên quyết} & Đã đăng nhập với tư cách xem/quản lý trạm. \\ \hline
    \textbf{Luồng sự kiện chính} &
      \begin{enumerate}[leftmargin=*, nosep, topsep=0pt]
        \item Người dùng mở màn hình quản lý trạm hạ cánh.
        \item Hệ thống hiển thị danh sách các trạm hạ cánh, sức chứa và trạng thái hoạt động.
        \item Nhập hoặc cập nhật thông tin trạm gồm tên trạm, địa chỉ, định vị lat/lng cùng công suất tối đa.
        \item Chọn trạng thái hoạt động mới: Đang hoạt động, Bảo trì hoặc Ngừng.
        \item Hệ thống lưu thay đổi và đồng bộ trạng thái trạm.
      \end{enumerate} \\ \hline
    \textbf{Luồng rẽ nhánh / Ngoại lệ} &
      \textbf{4a.} Trạm đang có đơn hàng chờ ở trạng thái Đang giao hoặc Đã lên lịch $\rightarrow$ Hệ thống cảnh báo không cho chuyển trạng thái sang Bảo trì/Ngừng trừ khi gán lại đơn hàng. \\ \hline
    \textbf{Hậu điều kiện}      & Thông tin trạm hạ cánh được cập nhật thành công. \\
    \hline
  \end{tabularx}
  \caption{Đặc tả Quản lý trạm hạ cánh}
  \label{tab:uc14}
\end{table}

% --- UC-15 ---
\subsection{UC-15: Xem báo cáo và thống kê}
\begin{table}[H]
  \centering
  \begin{tabularx}{\textwidth}{|l|X|}
    \hline
    \multicolumn{2}{|l|}{\textbf{Xem báo cáo và thống kê}} \\
    \hline
    \textbf{Tác nhân}           & Quản lý logistics, Quản trị viên \\ \hline
    \textbf{Mô tả}              & Cho phép hiển thị các báo cáo trực quan, thống kê hiệu suất giao hàng và hoạt động của Drone/Trạm. \\ \hline
    \textbf{Điều kiện tiên quyết} & Đã đăng nhập vào tài khoản quản trị/quản lý. \\ \hline
    \textbf{Luồng sự kiện chính} &
      \begin{enumerate}[leftmargin=*, nosep, topsep=0pt]
        \item Người dùng chọn mục ``Bảng điều khiển''.
        \item Hệ thống kết xuất dữ liệu và hiển thị các chỉ số cốt lõi bao gồm tổng đơn hàng, tỉ lệ thành công và hiệu suất Drone.
        \item Người dùng lọc dữ liệu theo khoảng thời gian ngày, tuần hoặc tháng.
        \item Hệ thống vẽ biểu đồ xu hướng giao hàng và thống kê sự cố.
      \end{enumerate} \\ \hline
    \textbf{Luồng rẽ nhánh / Ngoại lệ} &
      \textbf{2a.} Không tải được dữ liệu từ database $\rightarrow$ Báo lỗi và thử lại sau 3 giây. \\ \hline
    \textbf{Hậu điều kiện}      & Biểu đồ báo cáo và thống kê được kết xuất thành công. \\
    \hline
  \end{tabularx}
  \caption{Đặc tả Xem báo cáo và thống kê}
  \label{tab:uc15}
\end{table}

% --- UC-16 ---
\subsection{UC-16: Quản lý địa chỉ giao hàng}
\begin{table}[H]
  \centering
  \begin{tabularx}{\textwidth}{|l|X|}
    \hline
    \multicolumn{2}{|l|}{\textbf{Quản lý địa chỉ giao hàng}} \\
    \hline
    \textbf{Tác nhân}           & Khách hàng \\ \hline
    \textbf{Mô tả}              & Khách hàng lưu và cập nhật danh sách địa chỉ giao hàng nhận hàng quen thuộc. \\ \hline
    \textbf{Điều kiện tiên quyết} & Khách hàng đã đăng nhập. \\ \hline
    \textbf{Luồng sự kiện chính} &
      \begin{enumerate}[leftmargin=*, nosep, topsep=0pt]
        \item Khách hàng mở mục ``Sổ địa chỉ''.
        \item Chọn thêm mới địa chỉ, hoặc chọn địa chỉ cũ cần chỉnh sửa/xóa.
        \item Điền địa chỉ cụ thể, tên thành phố, định vị lat/lng tương ứng.
        \item Lọc và lưu địa chỉ mặc định.
      \end{enumerate} \\ \hline
    \textbf{Luồng rẽ nhánh / Ngoại lệ} &
      \textbf{3a.} Tọa độ lat/lng ngoài vùng phục vụ $\rightarrow$ Báo lỗi địa chỉ không hợp lệ. \\ \hline
    \textbf{Hậu điều kiện}      & Sổ địa chỉ của khách hàng được cập nhật trong cơ sở dữ liệu. \\
    \hline
  \end{tabularx}
  \caption{Đặc tả Quản lý địa chỉ giao hàng}
  \label{tab:uc16}
\end{table}

% --- UC-17 ---
\subsection{UC-17: Đăng ký tài khoản mới}
\begin{table}[H]
  \centering
  \begin{tabularx}{\textwidth}{|l|X|}
    \hline
    \multicolumn{2}{|l|}{\textbf{Đăng ký tài khoản mới}} \\
    \hline
    \textbf{Tác nhân}           & Khách hàng chưa đăng nhập \\ \hline
    \textbf{Mô tả}              & Người dùng tạo tài khoản Khách hàng mới để sử dụng hệ thống. \\ \hline
    \textbf{Điều kiện tiên quyết} & Chưa đăng nhập hệ thống. \\ \hline
    \textbf{Luồng sự kiện chính} &
      \begin{enumerate}[leftmargin=*, nosep, topsep=0pt]
        \item Khách hàng chọn nút ``Đăng ký''.
        \item Nhập họ và tên, email, mật khẩu và số điện thoại.
        \item Hệ thống kiểm tra tính duy nhất của email và độ mạnh của mật khẩu.
        \item Tạo mới tài khoản khách hàng, gửi mã kích hoạt hoặc chuyển sang đăng nhập.
      \end{enumerate} \\ \hline
    \textbf{Luồng rẽ nhánh / Ngoại lệ} &
      \textbf{3a.} Email đã tồn tại trong hệ thống $\rightarrow$ Báo lỗi trùng lặp. \\ \hline
    \textbf{Hậu điều kiện}      & Tạo tài khoản khách hàng mới thành công. \\
    \hline
  \end{tabularx}
  \caption{Đặc tả Đăng ký tài khoản mới}
  \label{tab:uc17}
\end{table}

% --- UC-18 ---
\subsection{UC-18: Quản lý thông tin cá nhân}
\begin{table}[H]
  \centering
  \begin{tabularx}{\textwidth}{|l|X|}
    \hline
    \multicolumn{2}{|l|}{\textbf{Quản lý thông tin cá nhân}} \\
    \hline
    \textbf{Tác nhân}           & Tất cả tác nhân \\ \hline
    \textbf{Mô tả}              & Mỗi người dùng cập nhật thông tin hồ sơ cá nhân và đổi mật khẩu. \\ \hline
    \textbf{Điều kiện tiên quyết} & Người dùng đã đăng nhập thành công. \\ \hline
    \textbf{Luồng sự kiện chính} &
      \begin{enumerate}[leftmargin=*, nosep, topsep=0pt]
        \item Người dùng chọn trang ``Thông tin cá nhân''.
        \item Chỉnh sửa họ tên, số điện thoại hoặc yêu cầu đổi mật khẩu.
        \item Hệ thống xác thực và cập nhật thông tin mới.
      \end{enumerate} \\ \hline
    \textbf{Luồng rẽ nhánh / Ngoại lệ} &
      \textbf{2a.} Mật khẩu cũ nhập sai $\rightarrow$ Báo lỗi mật khẩu cũ không đúng. \\ \hline
    \textbf{Hậu điều kiện}      & Hồ sơ thông tin cá nhân được cập nhật thành công. \\
    \hline
  \end{tabularx}
  \caption{Đặc tả Quản lý thông tin cá nhân}
  \label{tab:uc18}
\end{table}

% --- UC-19 ---
\subsection{UC-19: Tóm tắt thông tin giao hàng bằng AI}
\begin{table}[H]
  \centering
  \begin{tabularx}{\textwidth}{|l|X|}
    \hline
    \multicolumn{2}{|l|}{\textbf{Tóm tắt thông tin giao hàng bằng AI}} \\
    \hline
    \textbf{Tác nhân}           & Khách hàng, Điều phối viên, Quản lý logistics \\ \hline
    \textbf{Mô tả}              & Sử dụng LLM để tóm tắt lộ trình giao nhận, thông tin đơn hàng hoặc lịch sử sự cố thành văn bản báo cáo ngắn gọn. \\ \hline
    \textbf{Điều kiện tiên quyết} & Đơn hàng đã có dữ liệu giao hàng trong hệ thống. \\ \hline
    \textbf{Luồng sự kiện chính} &
      \begin{enumerate}[leftmargin=*, nosep, topsep=0pt]
        \item Người dùng bấm nút ``Tóm tắt bằng AI'' ở chi tiết chuyến giao hàng.
        \item Hệ thống lấy dữ liệu thô bao gồm thời gian bay, trạm đi/đến và trạng thái sự cố nếu có.
        \item Gửi dữ liệu qua API đến dịch vụ LLM để xử lý tóm tắt.
        \item Hiện đoạn tóm tắt giao hàng ngắn cho người dùng.
      \end{enumerate} \\ \hline
    \textbf{Luồng rẽ nhánh / Ngoại lệ} &
      \textbf{3a.} API ngoại vi bị lỗi $\rightarrow$ Báo lỗi không phản hồi, hiển thị tóm tắt thô từ hệ thống. \\ \hline
    \textbf{Hậu điều kiện}      & Hiển thị thành công đoạn văn bản tóm tắt hành trình giao hàng. \\
    \hline
  \end{tabularx}
  \caption{Đặc tả Tóm tắt thông tin giao hàng bằng AI}
  \label{tab:uc19}
\end{table}

% ============================================================
\section{Yêu Cầu Chức Năng}
% ============================================================

\subsection{Tổng Quan Chức Năng Hệ Thống}

Hệ thống SmartDroneDelivery được tổ chức thành bốn nhóm chức năng chính:

\begin{itemize}
  \item \textbf{Quản lý tài khoản và phân quyền:} Đăng nhập, đăng ký, quản lý người dùng theo vai trò.
  \item \textbf{Quản lý đơn hàng và kiện hàng:} Tạo, duyệt, lên lịch, theo dõi và xác nhận giao hàng.
  \item \textbf{Quản lý trạm hạ cánh:} Thêm, sửa, xóa trạm; cập nhật trạng thái hoạt động.
  \item \textbf{Dịch vụ AI hỗ trợ:} Ước lượng ETA và chatbot trợ lý khách hàng.
\end{itemize}

\subsection{Danh Sách Yêu Cầu Chức Năng}

\subsubsection{FR-01: Đăng Nhập Hệ Thống}
\textbf{Mô tả:} Cho phép tất cả người dùng đăng nhập vào hệ thống theo tài khoản được cấp.

\textbf{Chức năng:}
\begin{itemize}
  \item Nhập email/tên đăng nhập và mật khẩu.
  \item Xác thực tài khoản, cấp phiên làm việc.
  \item Phân quyền và điều hướng đến giao diện theo vai trò.
  \item Hỗ trợ đăng xuất khỏi hệ thống.
\end{itemize}

\subsubsection{FR-02: Quản Lý Tài Khoản Người Dùng}
\textbf{Mô tả:} Admin quản lý toàn bộ tài khoản trong hệ thống.

\textbf{Chức năng:}
\begin{itemize}
  \item Tạo tài khoản mới với vai trò tương ứng.
  \item Chỉnh sửa thông tin tài khoản.
  \item Khóa hoặc mở khóa tài khoản.
  \item Ghi log hoạt động quản trị.
\end{itemize}

\subsubsection{FR-03: Quản Lý Đơn Giao Hàng}
\textbf{Mô tả:} Cho phép Khách hàng tạo, theo dõi và quản lý đơn giao hàng; Dispatcher phê duyệt và lên lịch.

\textbf{Chức năng:}
\begin{itemize}
  \item Tạo đơn giao hàng mới với thông tin gói hàng và địa chỉ.
  \item Cập nhật hoặc hủy đơn khi đơn chưa được xử lý.
  \item Phê duyệt hoặc từ chối đơn.
  \item Lập lịch giao hàng, gán trạm hạ cánh.
  \item Xử lý đơn giao thất bại.
\end{itemize}

\subsubsection{FR-04: Quản Lý Kiện Hàng Tại Trạm}
\textbf{Mô tả:} Nhân viên trạm xác nhận và cập nhật trạng thái kiện hàng trong quá trình vận chuyển.

\textbf{Chức năng:}
\begin{itemize}
  \item Quét hoặc nhập mã đơn để xác nhận nhận kiện hàng.
  \item Cập nhật trạng thái kiện hàng như Đã nhận tại trạm hoặc Đang giao.
  \item Xác nhận giao hàng thành công.
\end{itemize}

\subsubsection{FR-05: Quản Lý Trạm Hạ Cánh}
\textbf{Mô tả:} Quản lý thông tin và trạng thái các trạm hạ cánh của drone trong hệ thống.

\textbf{Chức năng:}
\begin{itemize}
  \item Thêm trạm hạ cánh mới gồm tên, tọa độ và sức chứa.
  \item Cập nhật thông tin trạm.
  \item Xóa trạm.
  \item Cập nhật trạng thái hoạt động như Đang hoạt động, Bảo trì hoặc Ngừng.
\end{itemize}

\subsubsection{FR-06: Theo Dõi Giao Hàng Thời Gian Thực}
\textbf{Mô tả:} Cho phép Khách hàng, Dispatcher và Nhân viên trạm theo dõi vị trí và trạng thái giao hàng.

\textbf{Chức năng:}
\begin{itemize}
  \item Hiển thị bản đồ theo dõi vị trí drone theo thời gian thực.
  \item Cập nhật trạng thái đơn liên tục qua Supabase Realtime.
  \item Hiển thị ETA (thời gian dự kiến giao hàng).
\end{itemize}

\subsubsection{FR-07: Dịch Vụ Ước Lượng ETA}
\textbf{Mô tả:} Hệ thống tự động tính toán thời gian dự kiến giao hàng dựa trên tọa độ và tốc độ trung bình.

\textbf{Chức năng:}
\begin{itemize}
  \item Tính khoảng cách từ trạm đến địa chỉ giao hàng.
  \item Áp dụng công thức ước lượng ETA theo tốc độ trung bình drone.
  \item Cập nhật ETA theo vị trí thực tế khi đơn đang giao.
\end{itemize}

\subsubsection{FR-08: Chatbot Trợ Lý AI}
\textbf{Mô tả:} Chatbot hỗ trợ Khách hàng trả lời câu hỏi liên quan đến đơn hàng thông qua LLM API.

\textbf{Chức năng:}
\begin{itemize}
  \item Giao diện chat trực tiếp trong ứng dụng.
  \item Tích hợp LLM API của Groq để xử lý ngôn ngữ tự nhiên.
  \item Tra cứu thông tin đơn hàng ngữ cảnh để trả lời chính xác.
  \item Fallback: thông báo lỗi và gợi ý liên hệ Dispatcher khi API không khả dụng.
\end{itemize}

\subsubsection{FR-09: Quản lý thiết bị Drone}
\textbf{Mô tả:} Quản trị viên quản lý vòng đời thiết bị drone trong hệ thống để phục vụ việc lập lịch giao hàng.

\textbf{Chức năng:}
\begin{itemize}
  \item Thêm mới thiết bị drone vào hệ thống.
  \item Cập nhật thông tin drone bao gồm trạng thái hoạt động, dung lượng pin và ngày bảo trì.
  \item Xóa hoặc vô hiệu hóa thiết bị drone.
  \item Giám sát danh sách drone và mức năng lượng pin thời gian thực.
\end{itemize}

\subsubsection{FR-10: Thống kê và báo cáo giao hàng}
\textbf{Mô tả:} Cung cấp các thống kê trực quan về hoạt động giao hàng dành cho Logistics Manager.

\textbf{Chức năng:}
\begin{itemize}
  \item Biểu diễn biểu đồ thống kê đơn hàng theo thời gian.
  \item Thống kê tổng số chuyến bay thành công và tỷ lệ sự cố phát sinh.
  \item Xuất báo cáo hoạt động của drone và các trạm hạ cánh.
\end{itemize}

\subsubsection{FR-11: Quản lý thông tin cá nhân và Sổ địa chỉ}
\textbf{Mô tả:} Cho phép tất cả người dùng cập nhật hồ sơ cá nhân và quản lý địa chỉ nhận hàng của riêng họ.

\textbf{Chức năng:}
\begin{itemize}
  \item Chỉnh sửa hồ sơ cá nhân, họ tên, số điện thoại và đổi mật khẩu.
  \item Khách hàng thêm mới, cập nhật, xóa địa chỉ giao hàng trong sổ địa chỉ.
\end{itemize}

\subsubsection{FR-12: Tóm tắt thông tin giao hàng bằng AI}
\textbf{Mô tả:} Sử dụng trí tuệ nhân tạo để tóm tắt các chặng bay giao nhận và báo cáo sự cố tự động.

\textbf{Chức năng:}
\begin{itemize}
  \item Thu thập dữ liệu hành trình drone và các chặng giao hàng.
  \item Gọi API ngoại vi của LLM để sinh văn bản tóm tắt lộ trình.
\end{itemize}

\subsubsection{FR-13: Đăng ký tài khoản khách hàng}
\textbf{Mô tả:} Cho phép khách hàng vãng lai đăng ký tài khoản tự phục vụ trên hệ thống.

\textbf{Chức năng:}
\begin{itemize}
  \item Nhập thông tin đăng ký bao gồm email, mật khẩu, họ tên và số điện thoại.
  \item Xác thực định danh email và kích hoạt tài khoản sử dụng Supabase Auth.
\end{itemize}

% ============================================================
\section{Yêu Cầu Phi Chức Năng}
% ============================================================

\subsection{Xác Thực và Phân Quyền}
\begin{itemize}
  \item Hệ thống phải hỗ trợ cơ chế đăng nhập an toàn cho tất cả người dùng.
  \item Phân quyền dựa trên vai trò RBAC cho các nhóm người dùng: Khách hàng, Điều phối viên, Nhân viên vận hành trạm, Quản lý logistics, Quản trị viên hệ thống.
  \item Các thao tác nhạy cảm như xóa đơn hàng, thay đổi trạng thái giao hàng chỉ được thực hiện bởi người dùng có quyền hạn tương ứng.
\end{itemize}

Để làm rõ hơn phạm vi quyền hạn và các chức năng hệ thống phân phối cho từng nhóm người dùng, dưới đây là ma trận phân quyền sử dụng chức năng chi tiết:

\begin{table}[H]
  \centering
  \small
  \begin{tabularx}{\textwidth}{|l|X|c|c|c|c|c|}
    \hline
    \textbf{UC ID} & \textbf{Tên Use Case} & \textbf{C} & \textbf{SO} & \textbf{D} & \textbf{LM} & \textbf{A} \\ \hline
    UC-01 & Đăng nhập và Phân quyền                  & \checkmark & \checkmark & \checkmark & \checkmark & \checkmark \\ \hline
    UC-02 & Tạo đơn giao hàng mới                    & \checkmark &            &            &            &            \\ \hline
    UC-03 & Cập nhật và Hủy đơn giao hàng             & \checkmark &            &            &            &            \\ \hline
    UC-04 & Phê duyệt yêu cầu giao hàng              &            &            & \checkmark &            &            \\ \hline
    UC-05 & Lập lịch giao hàng                       &            &            & \checkmark &            &            \\ \hline
    UC-06 & Xử lý giao hàng thất bại                 &            &            & \checkmark &            &            \\ \hline
    UC-07 & Xác nhận nhận/gửi kiện hàng tại trạm     &            & \checkmark &            &            &            \\ \hline
    UC-08 & Theo dõi giao hàng thời gian thực        & \checkmark & \checkmark & \checkmark & \checkmark &            \\ \hline
    UC-09 & Xác nhận giao hàng thành công            & \checkmark & \checkmark &            &            &            \\ \hline
    UC-10 & Ước lượng thời gian giao hàng dự kiến    & \checkmark &            & \checkmark &            &            \\ \hline
    UC-11 & Trợ lý AI cho khách hàng                 & \checkmark &            &            &            &            \\ \hline
    UC-12 & Tạo và quản lý tài khoản người dùng      &            &            &            &            & \checkmark \\ \hline
    UC-13 & Quản lý thiết bị Drone                   &            &            &            &            & \checkmark \\ \hline
    UC-14 & Quản lý trạm hạ cánh                     &            &            &            & \checkmark & \checkmark \\ \hline
    UC-15 & Xem báo cáo và thống kê      &            &            &            & \checkmark & \checkmark \\ \hline
    UC-16 & Quản lý địa chỉ giao hàng                & \checkmark &            &            &            &            \\ \hline
    UC-17 & Đăng ký tài khoản mới                    & \checkmark &            &            &            &            \\ \hline
    UC-18 & Quản lý thông tin cá nhân                & \checkmark & \checkmark & \checkmark & \checkmark & \checkmark \\ \hline
    UC-19 & Tóm tắt thông tin giao hàng bằng AI      & \checkmark &            & \checkmark & \checkmark &            \\ \hline
  \end{tabularx}
  \begin{spacing}{0.95}
  \vspace{0.2cm}
  {\footnotesize \textit{Ghi chú ký hiệu viết tắt các tác nhân:} \\
  - \textbf{C}: Khách hàng \\
  - \textbf{SO}: Nhân viên vận hành trạm \\
  - \textbf{D}: Điều phối viên \\
  - \textbf{LM}: Quản lý logistics \\
  - \textbf{A}: Quản trị viên hệ thống \\
  - \textbf{\checkmark}: Biểu thị vai trò có quyền thao tác thực hiện Use Case.}
  \end{spacing}
  \caption{Ma trận phân quyền sử dụng chức năng hệ thống SmartDroneDelivery}
  \label{tab:permission-matrix}
\end{table}


\subsection{Thời Gian Phản Hồi}
\begin{itemize}
  \item Thời gian phản hồi cho các thao tác thông dụng như tạo đơn hoặc xem trạng thái phải dưới 3 giây.
  \item Các giao dịch quan trọng như xác nhận giao hàng hoặc cập nhật trạng thái drone phải hoàn tất trong vòng 10 giây.
\end{itemize}

\subsection{Hỗ Trợ Người Dùng}
\begin{itemize}
  \item Hệ thống phải hỗ trợ tối thiểu 100 người dùng đồng thời với hiệu suất ổn định và không làm giảm khả năng xử lý.
  \item Kiến trúc hệ thống phải cho phép mở rộng dễ dàng khi số lượng người dùng và đơn hàng tăng lên.
\end{itemize}

\subsection{RESTful API Phục Vụ Tích Hợp Bên Ngoài}
\begin{itemize}
  \item Cung cấp API chuẩn RESTful để tích hợp với các hệ thống logistics, kho vận hoặc dịch vụ bên ngoài.
  \item API phải có tài liệu hướng dẫn và đặc tả đầy đủ theo chuẩn Swagger hoặc OpenAPI.
\end{itemize}

\subsection{Triển Khai Bằng Supabase}
\begin{itemize}
  \item Hệ thống phải được triển khai và quản lý trên nền tảng Supabase để đảm bảo tính linh hoạt và dễ bảo trì.
  \item Supabase hỗ trợ cơ sở dữ liệu PostgreSQL, xác thực người dùng và API tích hợp, giúp giảm tải cấu hình thủ công.
\end{itemize}

\subsection{Sự Linh Hoạt}
\begin{itemize}
  \item Giao diện Web phải tương thích tốt với nhiều loại thiết bị và độ phân giải màn hình khác nhau để hiển thị tương thích trên cả máy tính, máy tính bảng và điện thoại di động.
  \item Thiết kế giao diện phải thân thiện, trực quan, dễ sử dụng cho mọi nhóm người dùng.
\end{itemize}

\subsection{Giám Sát Hệ Thống và Sao Lưu}
\begin{itemize}
  \item Hệ thống phải có cơ chế giám sát thời gian thực cho drone, trạm hạ cánh và tiến trình giao hàng.
  \item Hỗ trợ cơ chế sao lưu dữ liệu định kỳ và khôi phục nhanh chóng khi xảy ra sự cố.
\end{itemize}

\subsection{Kiến Trúc Phần Mềm và Khả Năng Bảo Trì}
\begin{itemize}
  \item Kiến trúc phần mềm phải được thiết kế theo dạng module hóa, tách biệt rõ ràng giữa các thành phần hệ thống.
  \item Đảm bảo mã nguồn dễ bảo trì, nâng cấp và mở rộng tính năng mới khi có nhu cầu trong tương lai.
\end{itemize}

% ============================================================
\section{Phụ Lục Yêu Cầu}
% ============================================================

\subsection{Quy Tắc Nghiệp Vụ}

\begin{longtable}{|l|p{12cm}|}
  \hline
  \textbf{Mã} & \textbf{Mô tả quy tắc} \\
  \hline
  \endhead
  \hline
  \endfoot
  BR-01 & Cân nặng gói hàng tối đa là 5kg; vượt quá không được tạo đơn. \\ \hline
  BR-02 & Khách hàng chỉ được hủy hoặc sửa đơn khi đơn ở trạng thái ``Chờ duyệt''. \\ \hline
  BR-03 & Mỗi đơn hàng chỉ được gán một trạm hạ cánh tại một thời điểm. \\ \hline
  BR-04 & Trạm đang trong trạng thái ``Bảo trì'' hoặc ``Ngừng'' không được gán đơn hàng mới. \\ \hline
  BR-05 & Mật khẩu tối thiểu 8 ký tự, gồm chữ hoa, chữ thường và số. \\ \hline
  BR-06 & Mỗi tài khoản chỉ được đăng nhập đồng thời trên một thiết bị với phiên duy nhất. \\ \hline
  BR-07 & Hệ thống tự động cập nhật ETA mỗi 30 giây khi đơn đang trong trạng thái ``Đang giao''. \\ \hline
  BR-08 & Sau 24 giờ không xác nhận giao hàng, hệ thống tự động chuyển sang ``Chờ xác nhận quá hạn''. \\
\end{longtable}

\subsection{Yêu Cầu Chung}

\begin{itemize}
  \item Hệ thống phải hoạt động ổn định trên nền tảng Web, hỗ trợ hiển thị responsive mượt mà trên các thiết bị di động sử dụng iOS và Android thông qua trình duyệt.
  \item Giao diện hỗ trợ đầy đủ tiếng Việt; các thông báo lỗi hiển thị rõ ràng bằng tiếng Việt.
  \item Dữ liệu nhạy cảm như mật khẩu và thông tin khách hàng phải được mã hóa và không lưu dưới dạng plain text.
  \item Mọi API endpoint phải yêu cầu xác thực JWT; tài nguyên không có token hợp lệ sẽ bị từ chối (\texttt{401 Unauthorized}).
  \item Hệ thống phải có cơ chế ghi log mọi thao tác quan trọng như tạo, sửa hoặc xóa đơn, thay đổi trạng thái kiện hàng, quản lý tài khoản.
  \item Thư thông báo được gửi qua email hoặc push notification khi trạng thái đơn hàng thay đổi.
  \item Hệ thống phải tuân thủ GDPR cơ bản: không chia sẻ dữ liệu cá nhân khách hàng cho bên thứ ba không được ủy quyền.
\end{itemize}

\subsection{Danh Sách Thông Báo Hệ Thống}

\begin{longtable}{|l|l|p{9cm}|}
  \hline
  \textbf{Mã} & \textbf{Loại} & \textbf{Nội dung thông báo} \\
  \hline
  \endfirsthead
  \hline
  \textbf{Mã} & \textbf{Loại} & \textbf{Nội dung thông báo} \\
  \hline
  \endhead
  \hline
  \endfoot
  \hline
  \endlastfoot
  MSG-001 & Thành công  & Đơn hàng đã được tạo thành công. \\ \hline
  MSG-002 & Thành công  & Đơn hàng đã được phê duyệt. \\ \hline
  MSG-003 & Thành công  & Đơn hàng đã được giao thành công. \\ \hline
  MSG-004 & Lỗi         & Sai tài khoản hoặc mật khẩu. Vui lòng thử lại. \\ \hline
  MSG-005 & Lỗi         & Tài khoản đã bị khóa. Vui lòng liên hệ quản trị viên. \\ \hline
  MSG-006 & Lỗi         & Gói hàng vượt tải trọng cho phép, tối đa là 5kg. \\ \hline
  MSG-007 & Lỗi         & Địa chỉ giao hàng nằm ngoài vùng phục vụ. \\ \hline
  MSG-008 & Lỗi         & Không có trạm hạ cánh khả dụng trong khu vực. \\ \hline
  MSG-009 & Cảnh báo    & Đơn hàng đã được xử lý, không thể thay đổi. \\ \hline
  MSG-010 & Cảnh báo    & Chatbot tạm thời không khả dụng. Vui lòng liên hệ điều phối viên. \\ \hline
  MSG-011 & Thông tin   & Phiên làm việc sắp hết hạn sau 5 phút. \\ \hline
  MSG-012 & Thông tin   & Đơn hàng của bạn đang được giao. thời gian dự kiến: \textit{[thời gian ước lượng]}. \\
\end{longtable}
